% ============================================================
%  CHAPTER 6 — VISUALIZATION: NOC DASHBOARD SYSTEM
% ============================================================
\chapter{Network Operations Center (NOC) Dashboard}

\section{Design Philosophy}

Traditional simulation reporting relies on monolithic, single-window dashboards that attempt to display all metrics simultaneously. This approach introduces two critical problems: (1)~UI rendering bottlenecks that steal CPU cycles from the simulation engine, and (2)~information overload that reduces operator situational awareness.

Aegis-IoT adopts a \textbf{decoupled, multi-window NOC architecture} where each dashboard executes as an independent Java Swing \texttt{JFrame} on its own thread pool. The dashboards communicate with the simulation core through thread-safe \texttt{SwingUtilities.invokeLater()} calls, ensuring zero contention with the discrete-event engine.

\section{Dashboard Overview}

\begin{table}[H]
\centering
\caption{NOC Dashboard Windows}
\label{tab:dashboards}
\begin{tabularx}{\textwidth}{@{}llX@{}}
\toprule
\textbf{Window} & \textbf{Class} & \textbf{Visualizations} \\
\midrule
Latency Hub & \texttt{LatencyDash} & Real-time actual vs.\ predicted latency line chart; t+10 autoregressive forecast curve \\
Security Radar & \texttt{SecurityDash} & 2D scatter plot (packet size vs.\ latency) with anomaly highlighting; OneClassSVM decision boundary \\
Telemetry Matrix & \texttt{TelemetryDash} & Live KPI counters (packets, anomalies, MAE); error distribution histogram \\
Energy \& Twin & \texttt{EnergyDash} & Battery depletion curve (100\% to 0\%); Digital Twin health radar with composite scoring \\
\bottomrule
\end{tabularx}
\end{table}

\section{Dashboard 1: Latency \& Forecasting Hub}

The Latency Hub displays two panels in a vertical split layout:

\subsection{Panel 1: Live Latency Tracking}

This panel renders a dual-line chart comparing actual flow durations (cyan) against AI-predicted values (magenta). The chart maintains a sliding window of the last 150 packets to prevent memory accumulation.

Key rendering details:
\begin{itemize}
    \item Anti-aliased rendering via \texttt{RenderingHints.VALUE\_ANTIALIAS\_ON}
    \item Dynamic Y-axis scaling: $\text{maxVal} \times 1.2$ for headroom
    \item Grid overlay with 5 horizontal reference lines
    \item 3px stroke width with round caps and joins for smooth curves
\end{itemize}

\subsection{Panel 2: Time-Series Forecasting}

This panel visualizes the autoregressive t+10 forecaster output in neon green. The forecasted values represent the model's prediction of network latency 10 packets into the future, enabling proactive congestion avoidance.

\section{Dashboard 2: Security \& Anomaly Radar}

The Security Radar provides a spatial visualization of the OneClassSVM decision landscape:

\begin{itemize}
    \item \textbf{Normal packets} are rendered as small cyan dots (4px radius, 120 alpha)
    \item \textbf{Anomalous packets} ($\text{score} < 0$) are rendered as large red circles (10px radius) with white outlines
    \item The X-axis maps packet size; the Y-axis maps flow duration
    \item A 5$\times$5 grid overlay provides spatial reference
\end{itemize}

This visualization allows operators to immediately identify DDoS clusters—groups of anomalous packets that cluster in the high-size, low-arrival-time region of the feature space.

\section{Dashboard 3: Telemetry \& Error Matrix}

The Telemetry Matrix provides aggregate system health metrics in a side-by-side layout:

\subsection{KPI Panel}

Three large-format counters display:
\begin{enumerate}
    \item \textbf{Total Packets Processed} (cyan) — Running count of all packets through the Cloud sink
    \item \textbf{GAN Anomalies Blocked} (red) — Count of packets flagged by the OneClassSVM
    \item \textbf{Global MAE} (magenta) — Running mean absolute error: $\text{MAE} = \frac{\sum|y_i - \hat{y}_i|}{n}$
\end{enumerate}

\subsection{Error Distribution Histogram}

The histogram bins the absolute prediction errors into 15 buckets and renders gradient-filled bars from neon magenta to transparent purple. This visualization reveals the error distribution shape—ideally concentrated near zero for a well-performing model.

\section{Dashboard 4: Energy Dynamics \& Digital Twin}

\subsection{Battery Depletion Curve}

The battery panel models IoT device energy consumption using a physics-inspired drain formula:

\begin{equation}
B_{t+1} = B_t - (0.005 + \text{packetSize} \times 0.00002)
\end{equation}

where $B_0 = 100\%$. The curve is rendered with color-coded segments:
\begin{itemize}
    \item $B > 50\%$: Neon green (healthy)
    \item $20\% < B \leq 50\%$: Orange (warning)
    \item $B \leq 20\%$: Red (critical)
\end{itemize}

The complete battery history is retained (up to 10,000 points) to display the full depletion trajectory from 100\% to 0\%.

\subsection{Digital Twin Health Radar}

The Digital Twin panel computes a composite health score:

\begin{equation}
H = B_{\text{current}} - (D_{\text{anomaly}} \times 50)
\end{equation}

where $D_{\text{anomaly}}$ is an exponentially weighted moving average of recent anomaly density:

\begin{equation}
D_{t+1} = 0.95 \times D_t + 0.05 \times \mathbb{1}[\text{anomalyScore} < 0]
\end{equation}

The health score is visualized as a pulsing concentric circle whose radius scales linearly with $H$. Two reference rings provide context, and the color transitions from cyan (healthy) through orange (degraded) to red (critical).

\section{Screen Layout}

All four dashboards are positioned in a 2$\times$2 quadrant layout optimized for 1920$\times$1080 displays:

\begin{table}[H]
\centering
\caption{Dashboard Screen Positions}
\label{tab:positions}
\begin{tabular}{@{}llll@{}}
\toprule
\textbf{Dashboard} & \textbf{Position (x, y)} & \textbf{Size} & \textbf{Quadrant} \\
\midrule
Latency Hub & (50, 50) & 700$\times$450 & Top-Left \\
Security Radar & (770, 50) & 700$\times$450 & Top-Right \\
Energy/Twin & (50, 520) & 700$\times$450 & Bottom-Left \\
Telemetry Matrix & (770, 520) & 700$\times$450 & Bottom-Right \\
\bottomrule
\end{tabular}
\end{table}

\section{Thread Safety}

All dashboard data updates are marshaled onto the Swing Event Dispatch Thread (EDT) using \texttt{SwingUtilities.invokeLater()}, preventing race conditions between the simulation engine and the UI rendering pipeline:

\begin{lstlisting}[language=Java, caption={Thread-Safe Dashboard Update}]
public void addData(double actual, double predicted) {
    javax.swing.SwingUtilities.invokeLater(() -> {
        actualLatencies.add(actual);
        predictedLatencies.add(predicted);
        if (actualLatencies.size() > 1000) {
            actualLatencies.remove(0);
            predictedLatencies.remove(0);
        }
        repaint();
    });
}
\end{lstlisting}
